%% Section 1: Introduction
\section{Introduction}
\label{sec:introduction}

The digital transformation of education has fundamentally reshaped how learners interact with knowledge, driven by the proliferation of mobile devices, cloud computing, and artificial intelligence~\cite{Crompton2017,Sharples2015}. Mobile learning (m-learning), defined as learning mediated through handheld devices such as smartphones and tablets, has emerged as a dominant paradigm owing to its ubiquity, portability, and capacity for situated, context-aware instruction~\cite{Kukulska2009,Chee2017}. By 2025, over six billion smartphone users worldwide have access to educational applications, creating unprecedented opportunities for anytime-anywhere learning. Yet despite this technological abundance, significant pedagogical challenges persist, particularly in the domain of vocabulary acquisition and reading comprehension~\cite{Nation2013,Schmitt2008}.

Vocabulary knowledge is widely recognized as the single strongest predictor of reading comprehension and academic achievement across disciplines~\cite{Stahl2005,Nation2013}. Students who encounter unfamiliar words in textbooks frequently experience a cascading comprehension failure: unknown vocabulary impedes sentence-level understanding, which in turn undermines paragraph coherence and, ultimately, the construction of meaningful domain knowledge~\cite{Laufer2010}. Traditional vocabulary instruction---relying on word lists, dictionary lookups, and rote memorization---has been shown to produce shallow, decontextualized knowledge that transfers poorly to authentic reading tasks~\cite{Schmitt2008,Webb2007}. What learners need instead is contextual, in-situ vocabulary support that preserves the relationship between words and the passages in which they appear~\cite{Hulstijn2001}.

Augmented Reality (AR) offers a compelling solution to this challenge by overlaying digital information directly onto the learner's physical environment~\cite{Azuma1997,Wu2013}. Unlike virtual reality, which replaces the real world entirely, AR supplements it---enabling learners to view explanatory content superimposed on the very textbook pages they are reading. Mobile AR, leveraging the built-in cameras and processing capabilities of smartphones, has demonstrated significant benefits in education, including enhanced spatial understanding, increased engagement, and improved knowledge retention~\cite{Akcayir2017,Ibanez2018,Garzon2019}. However, the majority of existing mobile AR educational applications function as content delivery systems, presenting static overlays without adapting to individual learner needs or tracking mastery over time~\cite{Bacca2014,Garzon2019}.

This lack of adaptivity represents a critical gap. Intelligent Tutoring Systems (ITS), which model learner knowledge and provide personalized instruction, have been shown to produce learning gains approaching those of one-on-one human tutoring~\cite{VanLehn2011,Anderson2001}. Techniques such as Bayesian Knowledge Tracing (BKT) enable fine-grained estimation of per-concept mastery probabilities, allowing the system to identify precisely which words a student has learned and which require further reinforcement~\cite{Corbett1994}. Similarly, scaffolding content according to Bloom's taxonomy---progressing from basic recall to application and analysis---ensures that explanations match the learner's cognitive readiness~\cite{Bloom1956,Krathwohl2002}. Yet ITS and AR have rarely been integrated within a single mobile platform, and even fewer systems incorporate gamification mechanics to sustain long-term engagement~\cite{Dey2018,Klopfer2012}.

Gamification---the application of game design elements such as points, badges, leaderboards, and streaks in non-game contexts---has demonstrated robust positive effects on student motivation, engagement, and persistence in educational settings~\cite{Deterding2011,Hamari2014,Sailer2017}. When aligned with established motivational frameworks such as Keller's ARCS model (Attention, Relevance, Confidence, Satisfaction)~\cite{Keller2010}, gamification can transform routine learning activities into intrinsically rewarding experiences. The combination of AR's immersive visuals, ITS's adaptive intelligence, and gamification's motivational scaffolding has the potential to create a learning environment that is simultaneously engaging, personalized, and pedagogically sound. However, to our knowledge, no existing mobile application integrates all four components---OCR text scanning, intelligent tutoring, AR overlays, and gamification---within a unified framework for vocabulary acquisition.

The design and development of such a complex, multi-component educational system demands rigorous modeling to ensure that pedagogical objectives are faithfully translated into technical specifications. The Unified Modeling Language (UML), the industry-standard notation for software architecture, provides the expressive power needed to capture both the structural organization (classes, relationships, actors) and the dynamic behavior (workflows, message sequences) of the system~\cite{Booch2005,Rumbaugh2004}. Recent work by Ouariach et al.~\cite{Ouariach2023,Ouariach2025,Ouariach2026} has demonstrated the effectiveness of UML-based pedagogical scenario modeling for e-learning systems, showing how class diagrams, use case diagrams, activity diagrams, and sequence diagrams can collectively bridge the gap between educational design and software implementation.

The present paper makes a twofold contribution. \textbf{First}, we propose a comprehensive UML modeling approach for AR-based educational systems, demonstrating how static views (class and use case diagrams) and dynamic views (activity and sequence diagrams) can be employed in concert to formalize the pedagogical scenario of a mobile AR vocabulary learning application. \textbf{Second}, we instantiate this approach in the design of AR-DeC (Augmented Reality Didactic Companion), a mobile application in which students scan textbook pages using their phone camera, OCR extracts individual words, an ITS engine generates adaptive explanations calibrated to the student's mastery level and Bloom's taxonomy stage, AR overlays present this content in the camera view, and a gamification engine awards points, badges, and streaks to sustain engagement. Through 15 UML classes organized into five functional groups, three actor roles, and detailed behavioral models, we provide a reusable architectural blueprint that can inform both the implementation of AR-DeC and the design of future mobile AR educational systems.

The remainder of this paper is organized as follows. \Cref{sec:theoretical} presents the theoretical framework, reviewing the literature on mobile AR in education, vocabulary acquisition, intelligent tutoring systems, gamification, and UML-based pedagogical modeling. \Cref{sec:modeling} constitutes the core contribution, detailing the UML modeling of the AR-DeC pedagogical scenario through class, use case, activity, and sequence diagrams. \Cref{sec:discussion} discusses the implications, contributions, and limitations of the proposed framework. Finally, \Cref{sec:conclusion} concludes the paper and outlines directions for future work.
